%=================================================================
% Section 6 — MIMO controller comparison (PLACEHOLDER — analysis not yet performed)
%=================================================================
\section{Toward a Combined MIMO Power--Drive Controller}
\label{sec:mimo}

% NOTE: This entire section is a placeholder. The analysis has NOT been performed.
% The purpose of including it in the plan is to (a) reserve the narrative slot,
% (b) define exactly what the comparison must contain so the future work is scoped,
% and (c) let reviewers of the outline see where the paper is headed.

The architecture presented so far treats the power-share loop (droop ratio $r$,
Section~\ref{sec:robust}) and the drive loop (motor current command from the velocity
controller) as decoupled SISO problems. In reality the two loops are coupled through the
bus: a drive transient changes the total load current, which both channels' droop
resistances convert into a share disturbance, while a share transient modulates the bus
voltage seen by the motor drive.
% NOTE: author to sharpen the physical coupling argument — the cross-coupling terms are
% visible in the full-order model (Section \ref{sec:robust}) as the bus-node interaction.

\subsection{Planned Comparison}
% TODO(analysis): none of the following exists yet. Planned scope:
\begin{itemize}
    \item Augment the share plant with the drive dynamics (VESC current loop approximated
          as a first-order lag + delay; motor/vehicle mechanical mode) to form a
          $2\times 2$ MIMO plant: inputs (droop ratio command, motor current command),
          outputs (measured share, wheel speed).
    \item Quantify the off-diagonal coupling (RGA / $\sigma$-plots) to establish whether
          the decentralized design is actually justified or merely convenient.
    \item Synthesize a MIMO $H_\infty$ controller on the same weight philosophy as
          Section~\ref{sec:robust} and compare against the decentralized pair on:
          worst-corner $\|S\|_\infty$, drive-transient share excursion, and regen-event
          bus excursion.
    \item Assess implementation cost on the Teensy 4.1 (state count vs.\ the current
          three-biquad realization at 1~kHz).
\end{itemize}

% FIGURE PLACEHOLDER
\begin{figure}[H]
    \centering
    \fbox{\parbox{0.9\linewidth}{\centering\vspace{2cm} TODO(figure): $2\times2$ MIMO block
    diagram — droop/share loop and drive loop with cross-coupling paths highlighted
    \vspace{2cm}}}
    \caption{Planned MIMO structure combining the power-share and drive loops.
    \textit{(Placeholder --- analysis not yet performed.)}}
    \label{fig:mimo_block}
\end{figure}

% PLANNING NOTES
% - Be explicit in the submitted paper that this is future work / outlook, OR perform the
%   RGA + synthesis before submission and promote it to a results section. Decide based on
%   page budget and thesis timeline.
% - The full-order 11-state model already contains the bus coupling; extending it with the
%   VESC lag is the cheapest path to the MIMO plant.
% - The EMS on the Pi commands share ≈ 1.0 during FC-charge cruise — the MIMO analysis
%   should include that operating point since the share loop is near an authority limit.
